\documentclass{article}
\usepackage{graphicx}
\usepackage[margin=1.5cm]{geometry}
\usepackage{amsmath}

\begin{document}
\twocolumn

\title{Chapter 1 In-Class Exercise: \texttt{print()} and Numeric Expressions}
\author{Prof. Jordan C. Hanson}

\maketitle

\section{Printing Arithmetic in Python}

\begin{enumerate}
\item Python can perform arithmetic directly inside the \texttt{print()} function.  The multiplication operator in Python is the asterisk, \texttt{*}.
\begin{itemize}

\item Create a new Python file called \texttt{multi\_table.py}.  Begin by printing a title:
\begin{verbatim}
print("Multiplication Table")
\end{verbatim}
\vspace{0.25cm}

\item We will arrange the products in rows and columns.  Begin by printing a row containing the column headings:
\begin{verbatim}
print("   ", 1, 2, 3, 4, 5)
\end{verbatim}
\vspace{0.25cm}

\item Now print the first row.  Python should calculate each product rather than having you type the answers:
\begin{verbatim}
print(1, 1*1, 1*2, 1*3, 1*4, 1*5)
\end{verbatim}
\vspace{0.25cm}

\item Run the program.  The numbers following the first \texttt{1} should be the products obtained by multiplying 1 by  1 through 5. \\ \vspace{0.25cm}

\end{itemize}
\end{enumerate}

\section{In-Class Exercise}

\begin{enumerate}

\item Complete the program using only simple Python statements, arithmetic expressions, and the \texttt{print()} function.  Your program should produce a two-dimensional multiplication table for the first five positive integers:

\begin{verbatim}
Multiplication Table

      1   2   3   4   5
1     1   2   3   4   5
2     2   4   6   8  10
3     3   6   9  12  15
4     4   8  12  16  20
5     5  10  15  20  25
\end{verbatim}

\begin{itemize}

\item Each row represents one integer from 1 through 5, and each column represents one integer from 1 through 5. \\ \vspace{0.25cm}

\item Each entry in the table should be calculated as the product of its row number and column number.  For example, the entry in row 3 and column 4 should be calculated using
\begin{verbatim}
3 * 4
\end{verbatim}
rather than by typing \texttt{12} directly. \\ \vspace{0.25cm}

\item Use one \texttt{print()} statement for each row of the multiplication table. \\ \vspace{0.25cm}

\item Before writing the entire program, determine the five multiplication expressions needed to produce row 4. Sketch this \texttt{print()} statement below before beginning in Google CoLab. \\ \vspace{0.75cm}

\item Run your program and compare its output with the example above.  Check that every row and column contains five products, and that the lower-right entry is the product of 5 and 5. \\ \vspace{0.5cm}

\end{itemize}
\end{enumerate}

\end{document}